\section{Pre-training versus tokenization}
\label{sec:decomp}

We now separate the two factors that a data-scaling curve alone cannot
distinguish: the value of pre-training and the value of tokenization. The
experiment is designed so that each can be isolated by changing exactly one
variable.

\subsection{Pre-training, at fixed tokenization, helps}

Holding tokenization fixed at non-overlapping 6-mer and data fixed at 50M
reads, we replace the pre-trained 29-layer NT-v2 backbone with a single-layer
randomly initialized transformer. The pre-trained backbone reaches 67.07\%
genus Top-1 versus 53.88\% for the random-initialized model---a $+13.19$ pp
pre-training advantage (Figure~\ref{fig:ablation}). Pre-training on genomic
sequence thus provides a real, transferable benefit for this task even though
the pre-training corpus is not microbial. This is the strongest case for GFMs
in metagenomics: the representation learned during pre-training is worth more
than 13 accuracy points that no amount of extra 6-mer data recovers.

\subsection{Tokenization, at fixed capacity, helps more}

Holding data fixed at 50M reads and training \emph{from scratch} (no
pre-training on either side), we vary only the tokenizer within the
MetaTransformer architecture. The effect dwarfs pre-training
(Table~\ref{tab:tok}): overlapping 13-mer reaches 87.42\% genus Top-1 while
overlapping 6-mer reaches only 48.87\% and non-overlapping 6-mer 47.00\%.
The overlapping-13-mer model trained from scratch (87.42\%) outperforms the
498M-parameter pre-trained NT-v2 with 6-mer tokenization (67.07\%) by
$+20.35$ pp on identical data. In other words, changing the tokenizer buys more
than changing from a shallow random network to a large pre-trained one. NT-v2's
non-overlapping 6-mer tokenization is the principal structural limitation of
the GFM approach here---not backbone capacity and not the adaptation strategy.

\begin{table*}[t]
\centering
\caption{\textbf{Tokenization governs both the ceiling and the ability to
scale.} Genus Top-1 accuracy (120 classes). MetaTransformer (MT) rows are
trained from scratch; NT-v2 is pre-trained. Only overlapping 13-mer reaches the
high-80s/90s and continues to improve from 50M to 250M reads, whereas 6-mer
representations plateau at 45--67\% regardless of pre-training or data.}
\label{tab:tok}
\begin{tabular*}{\textwidth}{@{\extracolsep{\fill}}lllcc@{}}
\toprule
Tokenizer & Overlap & Model & 50M & 250M \\
\midrule
13-mer & stride 1 (overlap)      & MT (scratch)      & 87.42\% & \textbf{98.7\%} \\
12-mer & stride 1 (overlap)      & MT (scratch)      & 78.83\% & --- \\
6-mer  & stride 1 (overlap)      & MT (scratch)      & 48.87\% & 50.6\% \\
6-mer  & non-overlap (stride 6)  & MT (scratch)      & 47.00\% & 45.0\% \\
6-mer  & non-overlap             & NT-v2 (pre-trained)& 67.07\% & 67.29\% \\
\botrule
\end{tabular*}
\end{table*}

Critically, tokenization also governs \emph{scalability}. As data grows from
50M to 250M, the overlapping-13-mer model improves by $+11$ pp
($87.5\%\!\rightarrow\!98.7\%$) while the 6-mer NT-v2 is flat ($+0.22$ pp;
Section~\ref{sec:scaling}). A tokenizer that preserves enough local signal both
raises the ceiling and keeps the model in the regime where more data helps.

\begin{figure}[t]
\centering
\includegraphics[width=\columnwidth]{backbone_ablation.pdf}
\caption{\textbf{Decomposing the sources of accuracy on a common 50M-read
baseline.} Left to right: a randomly initialized shallow 6-mer model; the
pre-trained NT-v2 6-mer model ($+13.2$ pp from pre-training, at fixed
tokenization); and an overlapping-13-mer model trained from scratch
($+20.4$ pp over pre-trained NT-v2, from tokenization alone). Tokenization is
the larger lever.}
\label{fig:ablation}
\end{figure}

\subsection{The ceiling is representational, not one of capacity or optimization}

A natural objection is that the 6-mer model is simply under-trained or
under-parameterized. The training-fit numbers refute this. The NT-v2 6-mer
model cannot even \emph{fit} the training set past $\sim$66\% (250M warm-start:
65.8\% train, 66.6\% validation---train and validation almost equal), whereas
the 13-mer model fits training to 98.9\% and validates at 97.5\%. When a model
cannot memorize the training data, the bottleneck is the input representation,
not capacity, data, or generalization: the non-overlapping 6-mer tokens simply
do not carry enough separable signal for a 150\,bp read.

\begin{figure}[t]
\centering
\includegraphics[width=\columnwidth]{train_fit.pdf}
\caption{\textbf{The ceiling is representational, not one of capacity.} Training
versus validation genus accuracy. The 6-mer model cannot fit the training set
past $\sim$66\% (train $\approx$ validation), whereas the overlapping-13-mer
model fits to 98.9\%. A model that cannot memorize its training data is limited
by its input representation, not by capacity or data.}
\label{fig:trainfit}
\end{figure}

\subsection{Is the 13-mer result just closed-set lookup?}
\label{sec:lookup}

The obvious worry about a 98.7\% number is that a 13-mer model is performing
disguised exact-match lookup against the reference---i.e.\ memorizing rather
than generalizing. We tested this directly by building non-neural $k$-mer
classifiers on the same 50M-read reference and the same test pool. A
Kraken-style unique-key lookup achieves only 0.72\% (because $k=13$ is
\emph{not} genus-unique: the $4^{13}=67.1$M 13-mer space is saturated by the
1{,}535 genomes, so test 13-mers are almost all ``seen'' but only $\sim$0.4\%
map to a single genus). A dominant-genus majority vote reaches 35.3\%, and a
full multinomial $k$-mer naive Bayes---using $P(k\text{-mer}\mid\text{genus})$
without any neural network---reaches 74.9\%. The MT 13-mer network adds a
further $+13$ pp over naive Bayes (87.5\%), capturing co-occurrence structure
beyond the naive-Bayes independence assumption. The mechanism is therefore
learned high-dimensional composition, \emph{not} exact-match lookup
(Figure~\ref{fig:kmer}).

\begin{figure}[t]
\centering
\includegraphics[width=\columnwidth]{kmer_baselines.pdf}
\caption{\textbf{Composition, not lookup.} Genus Top-1 accuracy of non-neural
$k$-mer classifiers on the same 50M-read reference and test pool, versus the
neural models. Exact-match unique-key lookup fails (0.7\%) because $k=13$ is not
genus-unique; a parameter-free multinomial 13-mer naive Bayes already reaches
74.9\%, \emph{exceeding} the 498M-parameter pre-trained NT-v2 (67\%). The neural
13-mer model adds a further $+13$ pp over naive Bayes.}
\label{fig:kmer}
\end{figure}

This analysis also delivers the paper's sharpest single comparison: a
parameter-free 13-mer naive Bayes (74.9\%) \emph{beats} the 498M-parameter
pre-trained NT-v2 with 6-mer tokenization (67\%). Representation ($k$-mer
length) is more decisive than model capacity or pre-training. We are explicit
about the accompanying caveat: because the reference genomes are nearly
completely covered by the test organisms, the 98.7\% figure characterizes
\emph{closed-set} classification against known genomes; novel-genome
generalization is a separate question that we flag as the primary limitation
(Section~\ref{sec:limitations}).
